\section{USGS 3DEP LiDAR data source and v2 train--test construction}
\label{sec:usgs-data-source}

The real spatial regression benchmark is constructed from the United States
Geological Survey (USGS) 3D Elevation Program (3DEP) LiDAR point cloud product.
We use the classified LiDAR point tiles from the USGS Rockyweb delivery
directory
\[
\texttt{https://rockyweb.usgs.gov/vdelivery/Datasets/Staged/Elevation/LPC/Projects/}
\]
under the project subdirectory
\[
\texttt{USGS\_LPC\_IL\_Winnebago\_2018/laz/}.
\]
The resulting regression task is two-dimensional ground elevation prediction:
the input is the horizontal coordinate pair \(x_i=(e_i,n_i)\), where \(e_i\)
and \(n_i\) are the projected easting and northing coordinates stored in the
LAS/LAZ files, and the response is the corresponding elevation \(z_i\).
Thus each sample is interpreted as
\[
(e_i,n_i) \longmapsto z_i .
\]

\paragraph{Download and tile-selection rules.}
For the v2 experimental process, the USGS datasets are described according to
the latest v2 preprocessing code, implemented in
\texttt{preprocess\_usgs\_3dep\_lidar\_contiguous\_v2.py}.  This script first
discovers all available LAS/LAZ files for the Winnebago project.  It checks the
standard Rockyweb manifest files
\texttt{0\_file\_download\_links.txt} and \texttt{file\_download\_links.txt};
if these manifests are unavailable or incomplete, it falls back to the Apache
directory index and recursively searches for files ending in \texttt{.las} or
\texttt{.laz}.  When a manifest contains stale cloud-storage links, only the
file basename is trusted and the URL is canonicalized back to the current
Rockyweb project directory.  This makes the data source reproducible even if
older manifest URLs have moved.

For the main size sweep up to \(3\times 10^7\) training samples, we select a
spatially contiguous block of \(64\) LiDAR tiles.  Tile coordinates are parsed
from filenames of the form
\[
\texttt{USGS\_LPC\_IL\_Winnebago\_2018\_a\_b.laz},
\]
where \(a\) and \(b\) are treated as grid coordinates.  The tiles are ordered by
their grid-coordinate ranks, and the algorithm chooses a compact near-square
block close to the center of the available project grid.  For the generated
datasets with \(10^5\), \(3\times 10^5\), \(10^6\), \(3\times 10^6\),
\(10^7\), and \(3\times 10^7\) training points, the selected block contains
tile coordinates from \((2559,2055)\) to \((2573,2069)\).  For the largest
streaming variants, the metadata records \texttt{tile\_selection=all}; the
\(3\times 10^8\) version uses \(438\) tiles, while the \(10^9\) version uses
\(999\) tiles.

\paragraph{Point filtering and v2 raw sample extraction.}
Each selected LAS/LAZ file is read with \texttt{laspy}.  By default, v2 keeps
only points with LAS classification label \(2\).  In the LAS classification
convention, class \(2\) denotes ground returns, so this filter removes
buildings, vegetation, and other non-ground objects before the elevation
regression problem is formed.  The metadata records \texttt{finite\_only=true},
and the subsequent bounding-box, scale, mean, and standard-deviation checks
guard against non-finite values in the retained arrays.  No per-file cap is
used in the recorded benchmark datasets.

The v2 code is designed for very large datasets and avoids materializing all
eligible points in memory.  It first scans the selected tiles and records the
number of eligible ground points in each file.  If the total number of eligible
points exceeds the requested point budget, the script allocates an exact number
of sampled points to each tile.  When supported by NumPy, this allocation uses a
multivariate hypergeometric draw, followed by uniform sampling without
replacement inside each tile; this is equivalent to drawing a fixed-size sample
from the union of all eligible ground points while keeping the implementation
streaming.  A deterministic proportional-allocation fallback is used on older
NumPy versions.  The random generator for the sampling allocation is
\[
\operatorname{rng}
=
\text{\texttt{np.random.default\_rng}}(\text{\texttt{seed}}),
\qquad \text{\texttt{seed}}=0.
\]
For a target training size \(N_{\mathrm{train}}\) and test fraction
\(\rho_{\mathrm{test}}\), the total number of retained clean points is
\[
N_{\mathrm{total}}
=
\left\lceil
\frac{N_{\mathrm{train}}}{1-\rho_{\mathrm{test}}}
\right\rceil .
\]
This choice ensures that, after the train--test construction below, the
training set has the requested target size.

\paragraph{V2 train--test construction.}
After filtering and optional fixed-size sampling, v2 constructs train and test
arrays directly on disk using NumPy memmaps.  Let
\[
\mathcal{D}=\{(x_i,z_i)\}_{i=1}^{N_{\mathrm{total}}}
\]
denote the retained clean point set.  The desired test size is
\[
N_{\mathrm{test}}
=
\operatorname{round}(\rho_{\mathrm{test}}N_{\mathrm{total}}),
\qquad
N_{\mathrm{train}}=N_{\mathrm{total}}-N_{\mathrm{test}}.
\]
Instead of forming a global permutation of all retained points, v2 allocates the
exact number of training samples to each tile from the already allocated sample
counts.  This second allocation uses
\[
\text{\texttt{np.random.default\_rng}}(\text{\texttt{seed}}+17).
\]
The remaining sampled points in each tile are assigned to the test set.  Within
each tile, v2 then draws a local random permutation using the main generator
with \(\text{\texttt{seed}}=0\), writes the selected training coordinates and
elevations to \texttt{x\_train\_raw.npy} and \texttt{y\_train\_raw.npy}, and
writes the selected test coordinates and elevations to
\texttt{x\_test\_raw.npy} and \texttt{y\_test\_raw.npy}.  The metadata records
this split as
\[
\text{\texttt{method = random\_per\_file\_exact\_allocation}}.
\]
This is the v2 experimental path for large USGS runs: it is random,
reproducible from the recorded seeds, and avoids the out-of-memory behavior of
the older concatenate-then-split implementation.

For comparison with the requested target size, the final array sizes satisfy
\[
|\mathcal{D}_{\mathrm{train}}|=N_{\mathrm{train}},
\qquad
|\mathcal{D}_{\mathrm{test}}|=N_{\mathrm{test}}.
\]
For the standard benchmark datasets, \(\rho_{\mathrm{test}}=0.2\), giving an
\(80/20\) train--test split.  For the recorded \(10^9\)-training-sample
streaming run, the metadata records \(\rho_{\mathrm{test}}=0.08\), which gives
\(10^9\) training samples and \(86{,}956{,}522\) test samples from
\(1{,}086{,}956{,}522\) retained points.

\paragraph{Normalization.}
The train--test split is performed on the raw physical coordinates and
elevations.  Normalization statistics are then estimated from the training set
only and applied to both training and test data.  In v2 these statistics and
the normalization pass are computed chunk-by-chunk over the memmap arrays.  If
\(x_i=(e_i,n_i)\), define
\[
x_{\min}=\min_{i\in\mathcal{I}_{\mathrm{train}}} x_i
\quad\text{componentwise},
\qquad
s_x = \max_{i\in\mathcal{I}_{\mathrm{train}}}
\max_{k\in\{1,2\}} (x_{ik}-x_{\min,k}).
\]
The normalized input is
\[
\widetilde{x}_i = \frac{x_i-x_{\min}}{s_x}.
\]
For the target, the training-set mean and standard deviation are
\[
\mu_z=\frac{1}{N_{\mathrm{train}}}
\sum_{i\in\mathcal{I}_{\mathrm{train}}} z_i,
\qquad
\sigma_z =
\left(
\frac{1}{N_{\mathrm{train}}}
\sum_{i\in\mathcal{I}_{\mathrm{train}}}(z_i-\mu_z)^2
\right)^{1/2},
\]
and the normalized response is
\[
\widetilde{z}_i=\frac{z_i-\mu_z}{\sigma_z}.
\]
Using only training-set statistics prevents information from the test responses
from entering the training procedure.

\paragraph{Recorded dataset sizes.}
The processed datasets are saved as NumPy \texttt{.npz} files with arrays
\texttt{x\_train}, \texttt{y\_train}, \texttt{x\_test}, and \texttt{y\_test},
and each file is accompanied by a JSON metadata file recording the source URL,
downloaded tiles, filtering rule, random seed, split fraction, normalization
statistics, serialization mode, and final shapes.  In v2 the arrays are stored
as \(\text{\texttt{float32}}\); for datasets with at least \(5\times 10^7\) retained
points, the script uses uncompressed \texttt{np.savez} storage instead of
\texttt{np.savez\_compressed} to avoid excessive compression cost.  The latest
experiment runner resolves dataset files by the family prefix
\[
\texttt{USGS\_LPC\_IL\_Winnebago\_2018\_ground\_elevation\_regression}
\]
and the recorded \(N_{\mathrm{train}}\) value, so the experiment consumes the
processed stems of the form
\[
\texttt{USGS\_LPC\_IL\_Winnebago\_2018\_ground\_elevation\_regression\_ntrain}N.
\]
The main recorded USGS size sweep contains
\[
N_{\mathrm{train}}
\in
\{10^5,\;3\times10^5,\;10^6,\;3\times10^6,\;10^7,\;3\times10^7,\;
10^8,\;3\times10^8,\;10^9\}.
\]
For the \(80/20\) datasets, the test sizes are \(N_{\mathrm{test}}
=0.25N_{\mathrm{train}}\), e.g. \(10^6\) training points and \(250{,}000\)
test points.  The \(10^9\) streaming run is the exception noted above, with an
\(8\%\) test fraction.
